\section{A complete spectral core in the orbit norm}
\label{sec:spectral}
For a single angular mode use the unitary radial variable
\begin{equation}
 u(t,\theta)=\frac{e^{in\theta}}{\sqrt{2\pi\cosh t}}Q(t),\qquad
 b_\alpha(u)=\norm Q_2+
 \int_{\mathbb R}(1+|t|)^\alpha|Q(t)|\dd t.
 \label{eq:fractional-norm}
\end{equation}
Write the last integral as $M_\alpha(Q)$. The Casimir minus $1/4$ becomes
\begin{equation}
 H_k=-\partial_t^2+(1/4-k^2)\sech^2t,\qquad
 k=|n|,\quad D(H_k)=H^2(\mathbb R).
\end{equation}
All scattering estimates below are at fixed $k$ unless stated otherwise.

\subsection{The radial estimate for absolute orbit integration}
The repulsive operator $H_0$ has no bound states and no zero resonance.
For the latter, a bounded solution of $H_0h=0$ has
$h'(\pm\infty)=0$ and $h'=O(e^{-2|t|})$ by integration of the equation.
Integration by parts then gives
$\int(|h'|^2+\tfrac14\sech^2t|h|^2)\dd t=0$, so $h=0$.

Choose Jost solutions $F_{0,+}\sim e^{ist}$ at $+\infty$ and
$F_{0,-}(t,s)=F_{0,+}(-t,s)$:
\begin{align}
 F_{0,+}(t,s)&=e^{ist}{}_2F_1\left(\tfrac12,\tfrac12;1-is;
                         \tfrac{1-\tanh t}{2}\right),\\
 T_0(s)&=\frac{\Gamma(1/2-is)^2}{\Gamma(1-is)\Gamma(-is)},
 \qquad |T_0(s)|^2=\tanh^2\pi s,\\
 e_{0,\pm}(t,s)&=(2\pi)^{-1/2}T_0(s)F_{0,\pm}(t,s).
\end{align}
The distorted Fourier transform is unitary with two channels and measure
$\dd s$, by~\cite[Theorem 2.3]{GPR}. Volterra estimates on the good
half-axis and the scattering relation on the other half-axis give
\begin{equation}
 |e_{0,\pm}(t,s)|\le C\min\{1,s(1+|t|)\}.
 \label{eq:threshold-bound}
\end{equation}
For example $|F_{0,+}|\le C$ on $t\ge0$ and
$|F_{0,+}|\le C(1+|t|)$ at low energy on $t\le0$;
$T_0(s)=O(s)$ yields the second bound. The identity
$T_0F_{0,+}=\overline{F_{0,-}}+R_0F_{0,-}$, with $|R_0|\le1$,
yields the uniform bound on the other half-axis.

Let $f=|u|$, which is rotation invariant. Its radial transform is $|Q|$,
and let $q_f$ be the sum of the squares of its two spectral coefficients.
For $0<\alpha\le1$, \eqref{eq:threshold-bound} gives
\begin{equation}
 q_f(s)\le C M_\alpha(Q)^2\min(s^{2\alpha},1),\qquad
 \int_0^\infty q_f(s)\dd s=\norm Q_2^2.
\end{equation}
The spherical subrepresentation of $L^2(X)$ has principal spectrum of
multiplicity two: the radial Casimir is purely absolutely continuous above
$1/4$, while a final discrete representation has no $K=SO(2)$ fixed vector.
Its unit spherical matrix coefficient is
$P_{-1/2+is}(\cosh\lambda)$, with group Schur weight~\cite{MY,DLMF}
\begin{equation}
 \KP(s)=\frac{4\pi^2}{s\tanh\pi s}=\rhoP(s)^{-1}.
\end{equation}
The radial and group multiplicity bases differ by a measurable unitary
matrix, leaving $q_f$ unchanged. Consequently
\begin{equation}
 \norm{A_f}_{L^2(G)}^2=\int_0^\infty\KP(s)q_f(s)^2\dd s
 \le C_\alpha b_\alpha(u)^4,\qquad \tfrac14<\alpha\le1.
 \label{eq:autocorrelation-bound}
\end{equation}
At zero the integrable power is $s^{4\alpha-2}$; above one use the
boundedness of $\KP$, the bound on $q_f$, and Parseval. The equality is
first applied to smooth compact spectral densities and then extended by
simultaneous $L^1(\dd s)$ and $L^2(\KP\dd s)$ approximation. This does
not exchange arbitrary distributional Schur kernels.

Combining \eqref{eq:orbit-autocorrelation} and
\eqref{eq:autocorrelation-bound} yields
\begin{equation}
 \norm{a_{u_Lu_Rv}}_2
 \le C_\alpha b_\alpha(u_L)b_\alpha(u_R)\norm v.
 \label{eq:product-orbit-bound}
\end{equation}
This constant is independent of the angular labels, since it uses the
radial operator of the absolute values. It proves the compact-test
estimate used in Lemma~\ref{lem:position-distribution}.

\subsection{All angular sectors, including bound states}
Set $\ell_k=k-1/2$ and $A_k=\partial_t+\ell_k\tanh t$. Then
\begin{equation}
 A_k^*A_k=H_k+\ell_k^2,\qquad
 A_kA_k^*=H_{k-1}+\ell_k^2.
\end{equation}
The second operator is strictly positive, with lower bound $1/4$ for
$k=1$ and $2k-2$ for $k\ge2$. Polar decomposition gives an isometry onto
$(\ker A_k)^\perp$,
\begin{equation}
 V_k=A_k^*(H_{k-1}+\ell_k^2)^{-1/2},\qquad
 \ker A_k=\operatorname{span}\{\sech^{\ell_k}t\}.
\end{equation}
It intertwines $H_{k-1}$ and $H_k$. Induction proves the full decomposition
\begin{equation}
 H_k\simeq H_0\oplus\bigoplus_{\Delta=1}^k\{-(\Delta-1/2)^2\}.
 \label{eq:darboux-spectrum}
\end{equation}
The bound states $b_{k,\Delta}$ are exponentially decreasing and follow
from the same intertwiners. Their Casimir is $\Delta(1-\Delta)$.
The no-bound-state hypothesis of~\cite{GPR} was used only for $H_0$.

The Jost solutions and scattering coefficients are
\begin{align}
 F_{k,+}(t,s)&=e^{ist}{}_2F_1\left(\tfrac12-k,\tfrac12+k;1-is;
                         \tfrac{1-\tanh t}{2}\right),\\
 T_k(s)&=T_0(s)(-1)^k\prod_{j=1}^k\frac{j-1/2-is}{j-1/2+is},\\
 R_k(s)&=-\frac{i(-1)^k}{\sinh\pi s}T_k(s).
 \label{eq:full-scattering}
\end{align}
Thus $T_k(s)=-i(-1)^k\pi s+O(s^2)$ and $R_k(0)=-1$.
With $W_k=F_{k,-}F'_{k,+}-F'_{k,-}F_{k,+}$,
$W_k(0)=2(-1)^{k+1}/\pi\ne0$, excluding a zero resonance.
The normalized modes $e_{k,\pm}=(2\pi)^{-1/2}T_kF_{k,\pm}$ are the
unitary continuous basis: applying $V_k$ to the preceding modes multiplies
them by the unit phases $\mp(\ell_k-is)/\sqrt{\ell_k^2+s^2}$.

\subsection{Spectral cutoff in a stronger position topology}
For $Q\in\Cc(\mathbb R)$ its continuous coefficients satisfy
\begin{equation}
 \begin{gathered}
 a_\pm(s)=\int\overline{e_{k,\pm}(t,s)}Q(t)\dd t
          =s\beta_\pm(s)\quad\text{near }0,\qquad \beta_\pm\in C^\infty,\\
 \sup_{s\ge1}s^L|\partial_s^ja_\pm(s)|<\infty
       \quad\text{for every }j,L.
 \end{gathered}
 \label{eq:spectral-endpoints}
\end{equation}
The zero follows from $T_k=O(s)$ on compact position sets. For high
energy, integrate $H_k^M$ onto $Q$ and use compact-position bounds on the
spectral derivatives of $e_{k,\pm}$.

These derivative bounds follow from a uniform Volterra estimate. On
$t\ge0$ write $F_{k,+}=e^{ist}m_+$; the kernel for $m_+$ is
$D_s(v)=\int_0^ve^{2isr}\dd r$. Since
$|\partial_s^jD_s(v)|\le2^jv^{j+1}/(j+1)$ and the potential has all
polynomial moments with an exponential tail, successive parameter
differentiation yields
\begin{equation}
 \sup_{s\ge0}|\partial_s^j(m_+(t,s)-1)|\le C_{k,j}e^{-2t}.
 \label{eq:jost-remainder}
\end{equation}
The other good half-axis is identical. Gamma-ratio asymptotics in
\eqref{eq:full-scattering} give bounded derivatives of $T_k,R_k$ on
$[0,\infty)$. On the bad half-axis the scattering relation implies that
derivatives of the full modes grow only polynomially in $|t|$.

Choose $\chi=1$ on $[0,1]$, $\chi=0$ on $[2,\infty)$, smooth in between.
The low-spectral part is
\begin{equation}
 L_\epsilon Q(t)=\sum_\pm\int_0^\infty
 e_{k,\pm}(t,s)\chi(s/\epsilon)a_\pm(s)\dd s.
\end{equation}
On each half-axis its plane-wave amplitudes have the form
$c_\epsilon(s)=s\chi(s/\epsilon)B(s)$ with the first two derivatives of
$B$ bounded. Hence
\begin{equation}
 c_\epsilon(0)=0,\quad \norm{c_\epsilon}_1=O(\epsilon^2),\quad
 |c_\epsilon'(0)|+\norm{c_\epsilon''}_1=O(1).
\end{equation}
Two integrations by parts retain the second endpoint term and give
\begin{equation}
 |L_\epsilon Q(t)|\le C\min(\epsilon^2,|t|^{-2})+
 C\epsilon^2e^{-2|t|}
 \le\frac{C'\epsilon^2}{(1+\epsilon|t|)^2}.
\end{equation}
For $0\le\alpha<1$,
\begin{equation}
 M_\alpha(L_\epsilon Q)=O(\epsilon^{1-\alpha}),\qquad
 \norm{L_\epsilon Q}_2=O(\epsilon^{3/2}).
 \label{eq:low-cutoff}
\end{equation}
This controls the linked large-position/low-energy region.

At high energy \eqref{eq:spectral-endpoints} permits arbitrarily many
integrations by parts without endpoint terms. The tail beyond a smooth
cutoff at $R$ is bounded by
$C_{L,M}R^{-L}(1+|t|)^{-M}+C_LR^{-L}e^{-2|t|}$.
Retaining all bound states, define
\begin{align}
 Q_{\epsilon,R}(t)
 &=\sum_{\Delta=1}^k\ip{b_{k,\Delta}}Q\,b_{k,\Delta}(t)\notag\\
 &\quad+\sum_\pm\int_0^\infty e_{k,\pm}(t,s)
 (1-\chi(s/\epsilon))\chi(s/R)a_\pm(s)\dd s.
\end{align}
For $R>2\epsilon$ its error is exactly the low plus high tail. Thus
\begin{equation}
 \norm{Q-Q_{\epsilon,R}}_2+M_\alpha(Q-Q_{\epsilon,R})\longrightarrow0,
 \quad \epsilon\downarrow0,\ R\uparrow\infty,\quad \alpha<1.
 \label{eq:weighted-spectral-core}
\end{equation}
We henceforth take $1/4<\alpha<1$, for example $\alpha=1/2$.

\begin{theorem}[Completeness of the orbit-test core]
Let $\cS_{\mathrm{all}}$ be the span of products of two fixed-angular-mode
compact smooth continuous-spectral packets or bound modes, with arbitrary
$v\in\mathcal K$. It is contained in $\cD_{\mathrm{orb}}$. Every
$\phi\in\Cc(Y;\mathcal K)$ has approximants $\psi_j\in\cS_{\mathrm{all}}$
such that
\begin{equation}
 \norm{\phi-\psi_j}_{L^2(Y;\mathcal K)}+
 \norm{a_{\phi-\psi_j}}_2\longrightarrow0.
\end{equation}
In particular, $\overline{B\cS_{\mathrm{all}}}=\cF(\mathcal K)$.
\label{thm:ambient-core}
\end{theorem}
\begin{proof}
On a common compact position set, approximate $\phi$ uniformly by finite
matter-rank, finite angular-Fourier, smooth radial tensor sums. Fixed radial
cutoffs retain compact support. Lemma~\ref{lem:position-distribution}
controls this orbit error. For each of the finitely many factors use
\eqref{eq:weighted-spectral-core} and \eqref{eq:product-orbit-bound}.
A diagonal choice gives the limit. If $k$ is orthogonal to
$B\cS_{\mathrm{all}}$, it is orthogonal to all $B\phi$ with $\phi$ compact
smooth. Equation \eqref{eq:position-pairing} gives $Rk=0$ as a
distribution; injectivity gives $k=0$.
\end{proof}

\begin{remark}[A hard interior cutoff is different]
Deleting a narrow band around $s_*>0$ need not converge in this topology.
An amplitude nonzero at $s_*$ gives a packet of scale
$\epsilon f(\epsilon t)e^{is_*t}$, whose $M_\alpha$ generally scales as
$\epsilon^{-\alpha}$. The zero in \eqref{eq:spectral-endpoints} was
essential. The actual threshold will be imposed after identifying the
physical spectral Hilbert space.
\end{remark}
